\subsection*{Models}
\label{sec:axiom-systems-models}

\begin{tcolorbox}[colback=structurebox, colframe=structureborder, arc=2pt,
  left=8pt, right=8pt, top=6pt, bottom=6pt,
  title={\small\textbf{Model Presentation: Arithmetic as a Model}},
  fonttitle=\small\bfseries]
Reuse the arithmetic signature and arithmetic-style axioms from the previous
sections. A model begins by interpreting the symbols of that signature in an
actual mathematical structure.

\[
\begin{array}{rcl}
\text{Domain} & = & \text{the natural numbers},\\
0 & \mapsto & \text{the number zero},\\
S & \mapsto & \text{the successor operation},\\
+ & \mapsto & \text{addition}.
\end{array}
\]

In this structure, the statement
\[
  \forall x\;(x+0=x)
\]
holds because adding zero to any natural number gives back that same natural
number.

The formal symbol $0$ is only a symbol until it is interpreted as the number
zero. The formal symbol $+$ is only a symbol until it is interpreted as
addition. The symbols acquire meaning only after interpretation.
\end{tcolorbox}

\begin{remark*}[Structural remarks]
\textbf{Syntax vs meaning.}
Signatures, axioms, and theories are formal objects. They belong to syntax:
they specify vocabulary, adopted assertions, and consequences. Models provide
semantic meaning by assigning mathematical content to the symbols and by
determining which statements hold.

\textbf{One theory, many models.}
The same theory may have multiple models. Euclidean geometry may be realized
in different coordinate systems. The group axioms have many models, including
different concrete groups. The theory fixes the assertions; a model is one
mathematical realization in which those assertions hold.

\textbf{Countermodels.}
A statement fails to follow from a theory if there is a model of the theory in
which that statement is false. Countermodels are therefore evidence that a
claim is not forced by the axioms.

\textbf{Forward roadmap.}
The next sections use models to explain consistency and independence:
\[
\begin{array}{c}
\text{Models}\\
\Downarrow\\
\text{Consistency}\\
\Downarrow\\
\text{Independence}
\end{array}
\]
\end{remark*}

\begin{definitionbox}{Definition (Interpretation)}
\begin{definition}[Interpretation]
\label{def:axiom-systems-interpretation}
Let $\Sigma=(\mathcal{C},\mathcal{F},\mathcal{R},\operatorname{ar})$ be a
first-order signature, and let $M$ be a nonempty set. An
\emph{interpretation of $\Sigma$ on $M$} is an assignment $\mathcal{I}$ that
maps each symbol of $\Sigma$ to concrete data over $M$: constants are assigned
elements of $M$, function symbols are assigned operations on $M$, and relation
symbols are assigned relations on $M$.
\end{definition}
\end{definitionbox}

\begin{remark*}[Standard quantified statement]
For every $c\in\mathcal{C}$, $\mathcal{I}(c)\in M$. For every
$f\in\mathcal{F}$, if $\operatorname{ar}(f)=n$, then
\[
  \mathcal{I}(f):M^n\to M.
\]
For every $R\in\mathcal{R}$, if $\operatorname{ar}(R)=n$, then
\[
  \mathcal{I}(R)\subseteq M^n.
\]
\end{remark*}

\begin{remark*}[Interpretation]
An interpretation turns formal vocabulary into mathematical content. It tells
what each constant, function symbol, and relation symbol means in a particular
setting.
\end{remark*}

\begin{dependencies}
\begin{itemize}
  \item \hyperref[def:signature]{Signature}
\end{itemize}
\end{dependencies}

\begin{definitionbox}{Definition ($\mathcal{L}$-Structure)}
\begin{definition}[$\mathcal{L}$-Structure]
\label{def:axiom-systems-structure}
Let $\mathcal{L}$ be a first-order signature. An
\emph{$\mathcal{L}$-structure} is a pair
\[
  \mathfrak{M}=(M,\mathcal{I}),
\]
where $M$ is a nonempty domain and $\mathcal{I}$ is an interpretation of the
symbols of $\mathcal{L}$ on $M$.
\end{definition}
\end{definitionbox}

\begin{remark*}[Standard quantified statement]
For every first-order signature $\mathcal{L}$, an $\mathcal{L}$-structure
consists of a nonempty domain $M$ and an interpretation $\mathcal{I}$ such
that every constant symbol is assigned an element of $M$, every $n$-ary
function symbol is assigned a function $M^n\to M$, and every $n$-ary relation
symbol is assigned a relation on $M^n$.
\end{remark*}

\begin{remark*}[Interpretation]
A structure is the environment in which formal symbols receive meaning. The
domain supplies the objects, and the interpretation explains how the symbols
refer to objects, operations, and relations on that domain.
\end{remark*}

\begin{remark*}[Exposition]
In dependent type theory, this bundled data can be viewed schematically as
\[
\sum_{M:\mathrm{Type}}
\left(\prod_{c\in\mathcal{C}} M\right)
\times
\left(\prod_{f\in\mathcal{F}}(M^{\operatorname{ar}(f)}\to M)\right)
\times
\left(\prod_{R\in\mathcal{R}}(M^{\operatorname{ar}(R)}\to \mathrm{Prop})\right).
\]
This is the same semantic package represented by a Lean structure carrying a
carrier type, interpretations of constants and functions, and predicates for
relations.
\end{remark*}

\begin{dependencies}
\begin{itemize}
  \item \hyperref[def:signature]{Signature}
  \item \hyperref[def:axiom-systems-interpretation]{Interpretation}
\end{itemize}
\end{dependencies}

\begin{definitionbox}{Definition (Model)}
\begin{definition}[Model]
\label{def:axiom-systems-model}
Let $\mathcal{L}$ be a first-order signature, let $T$ be a theory over
$\mathcal{L}$, and let $\mathfrak{M}=(M,\mathcal{I})$ be an
$\mathcal{L}$-structure. The structure $\mathfrak{M}$ is a \emph{model of
$T$}, written $\mathfrak{M}\vDash T$, if every sentence in $T$ is true in
$\mathfrak{M}$.
\end{definition}
\end{definitionbox}

\begin{remark*}[Standard quantified statement]
For every theory $T$ over $\mathcal{L}$ and every $\mathcal{L}$-structure
$\mathfrak{M}$,
\[
  \mathfrak{M}\vDash T
  \quad\Longleftrightarrow\quad
  \forall \varphi\in T,\; \mathfrak{M}\vDash \varphi.
\]
\end{remark*}

\begin{remark*}[Interpretation]
A model is a structure that makes the theory true. It is not just an
interpretation of the vocabulary; it is an interpretation in which the adopted
assertions and their consequences hold.
\end{remark*}

\begin{dependencies}
\begin{itemize}
  \item \hyperref[def:axiom-systems-structure]{Structure}
  \item \hyperref[def:axiom-systems-theory]{Theory}
\end{itemize}
\end{dependencies}

\begin{definitionbox}{Definition (Satisfaction)}
\begin{definition}[Satisfaction]
\label{def:axiom-systems-satisfaction}
Let $\mathfrak{M}=(M,\mathcal{I})$ be an $\mathcal{L}$-structure. A
\emph{variable assignment} is a function $s:\operatorname{Var}\to M$. The
assignment extends recursively to a term-evaluation map $\bar{s}$ by
\[
  \bar{s}(x)=s(x),\qquad
  \bar{s}(c)=\mathcal{I}(c),
\]
and
\[
  \bar{s}\bigl(f(t_1,\ldots,t_n)\bigr)
  =
  \mathcal{I}(f)\bigl(\bar{s}(t_1),\ldots,\bar{s}(t_n)\bigr).
\]
The satisfaction relation $\mathfrak{M}\vDash\varphi[s]$ is then defined by
induction on formulas: atomic relations are tested inside the interpreted
relation, equality is interpreted as equality in $M$, Boolean connectives are
interpreted by the corresponding truth operations, and quantifiers range over
all elements of $M$.
\end{definition}
\end{definitionbox}
\begin{remark*}[Standard quantified statement]
For terms $t_1,\ldots,t_n$ and an $n$-ary relation symbol $R$,
\[
  \mathfrak{M}\vDash R(t_1,\ldots,t_n)[s]
  \quad\Longleftrightarrow\quad
  \bigl(\bar{s}(t_1),\ldots,\bar{s}(t_n)\bigr)\in\mathcal{I}(R).
\]
For equality,
\[
  \mathfrak{M}\vDash (t_1=t_2)[s]
  \quad\Longleftrightarrow\quad
  \bar{s}(t_1)=\bar{s}(t_2).
\]
For quantifiers,
\[
  \mathfrak{M}\vDash(\forall x\,\varphi)[s]
  \quad\Longleftrightarrow\quad
  \forall a\in M,\; \mathfrak{M}\vDash\varphi[s[x\mapsto a]],
\]
and
\[
  \mathfrak{M}\vDash(\exists x\,\varphi)[s]
  \quad\Longleftrightarrow\quad
  \exists a\in M,\; \mathfrak{M}\vDash\varphi[s[x\mapsto a]].
\]
\end{remark*}

\begin{remark*}[Predicate reading]
The predicate form records the relation or property named by the statement in
the displayed logical context.
\end{remark*}


\begin{remark*}[Interpretation]
Satisfaction is the semantic test for a statement. Once the symbols have been
interpreted in a structure, the statement can be checked as true or false in
that setting.
\end{remark*}

\begin{remark*}[Exposition]
If $\varphi$ is a sentence, then it has no free variables, so its truth does
not depend on the initial assignment $s$. In that case the notation is
shortened from $\mathfrak{M}\vDash\varphi[s]$ to $\mathfrak{M}\vDash\varphi$.
This is the bridge from satisfaction of a single sentence to the definition of
a model of a theory.
\end{remark*}

\begin{dependencies}
\begin{itemize}
  \item \hyperref[def:axiom-systems-structure]{Structure}
\end{itemize}
\end{dependencies}

\begin{definitionbox}{Definition (Countermodel)}
\begin{definition}[Countermodel]
\label{def:axiom-systems-countermodel}
Let $\Gamma\subseteq\operatorname{Sent}(\mathcal{L})$ be a set of premises,
let $\varphi\in\operatorname{Sent}(\mathcal{L})$, and let $\mathfrak{M}$ be
an $\mathcal{L}$-structure. A \emph{countermodel to the entailment
$\Gamma\vDash\varphi$} is a structure $\mathfrak{M}$ that satisfies every
sentence in $\Gamma$ but does not satisfy $\varphi$.
\end{definition}
\end{definitionbox}
\begin{remark*}[Standard quantified statement]
For $\Gamma\subseteq\operatorname{Sent}(\mathcal{L})$ and
$\varphi\in\operatorname{Sent}(\mathcal{L})$,
\[
  \mathfrak{M}\in\mathsf{CounterModel}(\Gamma,\varphi)
  \quad\Longleftrightarrow\quad
  \bigl(\forall \gamma\in\Gamma,\; \mathfrak{M}\vDash\gamma\bigr)
  \land
  \bigl(\mathfrak{M}\not\vDash\varphi\bigr).
\]
Equivalently, for classical closed sentences,
\[
  \mathfrak{M}\in\mathsf{CounterModel}(\Gamma,\varphi)
  \quad\Longleftrightarrow\quad
  \bigl(\forall \gamma\in\Gamma,\; \mathfrak{M}\vDash\gamma\bigr)
  \land
  \bigl(\mathfrak{M}\vDash\neg\varphi\bigr).
\]
\end{remark*}

\begin{remark*}[Predicate reading]
The predicate form records the relation or property named by the statement in
the displayed logical context.
\end{remark*}


\begin{remark*}[Interpretation]
A countermodel shows that the theory does not force the statement. It keeps
the axioms true while making the proposed statement false.
\end{remark*}

\begin{remark*}[Exposition]
If $\Gamma=\emptyset$, then a countermodel to validity of $\varphi$ is a
structure $\mathfrak{M}$ such that $\mathfrak{M}\vDash\neg\varphi$. For a
theory $T$, models of $T$ satisfying $\varphi$ and models of $T$ satisfying
$\neg\varphi$ witness the independence pattern: the theory alone does not
settle $\varphi$.
\end{remark*}

\begin{remark*}[Exposition]
In a mechanized setting, invalidity is witnessed by constructing semantic
data:
\[
  \exists\mathfrak{M}\,
  \bigl((\forall\gamma\in\Gamma,\;\mathfrak{M}\vDash\gamma)
  \land
  \mathfrak{M}\vDash\neg\varphi\bigr).
\]
The constructed structure and satisfaction proofs are the certificate that the
entailment cannot be valid.
\end{remark*}

\begin{dependencies}
\begin{itemize}
  \item \hyperref[def:axiom-systems-model]{Model}
  \item \hyperref[def:axiom-systems-satisfaction]{Satisfaction}
\end{itemize}
\end{dependencies}
